\chapter{\hybrid}
\label{ch:hybrid}
Alongside reproducing \tfunet{}, I designed a second network, which I call
\hybrid{}. It keeps the same overall U-Net shape but combines several
components that other RFI and segmentation work has used, chosen with the
properties of RFI in mind. This chapter describes the model and its results;
\cref{ch:drivers} then tests which of its components actually matter.

\section{Design}
\label{sec:hybrid_design}
The starting point was an observation about the synthetic data: after
normalisation, faint RFI sits only about 1.5 noise standard deviations above
the background. At that level a single pixel cannot be told apart from a
noise fluctuation. A faint narrowband line is detectable only because it
extends over many time samples, not because any one pixel is bright. Most of
the design choices were aimed at using this kind of spatial context.

\Cref{tab:hybrid_vs_tfunet} lists how the model differs from \tfunet{}.

\begin{table}[H]
  \centering
  \caption{Differences between \tfunet{} and \hybrid{}.}
  \label{tab:hybrid_vs_tfunet}
  \begin{tabular}{p{3.4cm}p{4.3cm}p{5.2cm}}
    \toprule
    Component & \tfunet{} & \hybrid{} \\
    \midrule
    Convolution blocks & two 3$\times$3 convolutions & residual blocks (with a shortcut path) \\
    Line-shaped filters & none & multi-scale strip convolutions \\
    Channel attention & none & efficient channel attention (ECA) \\
    Normalisation layers & none & GroupNorm after each convolution \\
    Output layer & ReLU, then softmax & raw scores, then softmax \\
    Padding & valid (output smaller than input) & same (output = input size) \\
    Depth / first-layer width & 3 levels / 32 & 4 levels / 8 \\
    Loss & cross-entropy & cross-entropy + Dice \\
    Framework & TensorFlow & PyTorch \\
    Parameters & 465{,}986 & 593{,}842 \\
    \bottomrule
  \end{tabular}
\end{table}

\paragraph{Residual blocks.} Each block adds its input back to its output
through a shortcut\footcite{he2016}, so a feature only has to learn a correction
rather than be rebuilt at every layer. The aim was to stop weak, low-contrast
features being washed out as they pass through many layers. RFI-Net, the best
method in Mesarcik et al.'s comparison, is also built from residual
blocks\footcite{yang2020rfinet}.

\paragraph{Multi-scale strip convolutions.} Ordinary filters are small
squares, but much RFI is line-shaped: narrowband RFI is a thin line along the
time axis, and a wideband burst a thin line along the frequency axis. A strip
filter of size $1\times K$ or $K\times1$ adds up the signal along such a line.
Averaging $K$ pixels of a coherent line improves its signal-to-noise ratio by
about $\sqrt{K}$, while noise averages down, which is how a 1.5$\sigma$ line
could become detectable. Filters of three lengths ($K = 7, 11, 21$) run in
parallel to cover RFI of different durations and bandwidths. Multi-scale
convolutions of this kind have been used for RFI before\footcite{gu2024emsca}, and
a closely related strip-convolution design has since been
published\footcite{mars2026}, so this component is not new.

\paragraph{Efficient channel attention.} After the strip filters, some feature
channels carry evidence of faint RFI and others of bright RFI. ECA\footcite{wang2020eca}
learns a weight for each channel at very little cost (a single small
one-dimensional convolution), so that the faint-structure channels are not
drowned out by the bright ones.

\paragraph{GroupNorm.} Normalisation layers rescale the values inside the
network so that they stay in a stable range. A common choice, batch
normalisation, estimates its statistics across the images in a batch, which is
unreliable when, as here, the GPU memory allows only one 512$\times$512 image
per batch. GroupNorm\footcite{wu2018groupnorm} computes its statistics within a
single image and so works at any batch size. \tfunet{} has no normalisation
layers at all.

\paragraph{No ReLU on the output.} \hybrid{} passes its raw output scores
directly to the softmax. This removes the failure described in
\cref{sec:relu_trap}, where class weighting and an output ReLU freeze
\tfunet{}.

\paragraph{Same padding.} The output has the same size as the input, so every
pixel of the 512$\times$512 image receives a prediction. With \tfunet{}'s valid
padding, a border around each image is lost.

\section{Training}
\label{sec:hybrid_training}
The model is trained with the sum of two losses. Cross-entropy scores each
pixel separately. The Dice loss\footcite{milletari2016} measures how well the
predicted RFI region overlaps the true one as a whole, which helps when the
positive class is rare, as RFI is in LOFAR. I used the Adam optimiser with a
learning rate of $10^{-3}$ and one image per batch, for 40 epochs of 700 steps
(28{,}000 gradient steps). As for \tfunet{}, the input is scaled with a fixed
range measured on training images only, and the best checkpoint and the
decision threshold are chosen on the 735 validation images.

\section{Results on synthetic data}
On the synthetic data the model performs very well, but as
\cref{sec:gap} showed, this data is much easier than real observations, so
these results are mainly a check that the model trains correctly.
On the 276$\times$600 dataset a wide version of the model (32 filters in the
first layer, 9{,}304{,}186 parameters) reached an \Fone{} of 0.9808, and on the
1024$\times$265 dataset the narrow version used in the rest of this report (8
filters, 593{,}842 parameters) reached 0.9878. These values use a threshold
chosen on the validation set.

The width of the network can be reduced a long way at little cost
(\cref{tab:width}). Going from 32 to 8 filters in the first layer removes 94\%
of the parameters and lowers \Fone{} by only 0.006; only at 4 filters does the
model clearly get worse. Each width was trained once, so differences of less
than about 0.01 should not be over-interpreted.

\begin{table}[H]
  \centering
  \caption{Effect of network width on the 276$\times$600 synthetic dataset
    (one run each).}
  \label{tab:width}
  \begin{tabular}{rrr}
    \toprule
    First-layer filters & Parameters & \Fone{} \\
    \midrule
    4  & 152{,}582     & 0.9547 \\
    8  & 593{,}842     & 0.9749 \\
    16 & 2{,}342{,}474 & 0.9788 \\
    32 & 9{,}304{,}186 & 0.9812 \\
    \bottomrule
  \end{tabular}
\end{table}

\section{Results on LOFAR}
\label{sec:hybrid_lofar}

\subsection{Class weighting}
I first trained the model on LOFAR with class weighting, giving RFI pixels 57
times the weight of clean ones to balance their rarity. Over three seeds it
reached a maximum \Fone{} of $0.6467 \pm 0.0079$. However, the threshold chosen
on the validation set was 0.96, far from the usual 0.5: the weighting pushed
the model's outputs so high that only near-certain predictions could be
trusted. Without class weighting, the score rose to $0.6603 \pm 0.0040$ and the
threshold fell to $0.26$. Removing the weight improved the result for every one
of the three seeds (by 0.009 to 0.016; paired $t = 6.03$ with 2 degrees of
freedom, significant at the 5\% level), so the unweighted model is used from
here on.

Interestingly, the area under the ROC curve moved the opposite way, falling
from 0.980 to 0.938. The ROC curve measures how well the model ranks RFI pixels
above clean ones, while \Fone{} measures the quality of the final yes/no
decision. Class weighting improved the ranking but made the decisions worse,
which shows that a high ROC area alone does not mean a model flags RFI well.

\subsection{Comparison with published methods}
\Cref{tab:final} and \cref{fig:final} compare the final model with the other
methods on the 109 human-labelled test images. \hybrid{} reaches a maximum
\Fone{} of $0.6603 \pm 0.0040$, which is 0.062 above RFI-Net (0.5979), the best
result reported by Mesarcik et al.; this difference is significant
(one-sample $t = 27.2$ with 2 degrees of freedom). With the threshold fixed in
advance on the validation set, the pooled \Fone{} is $0.6561 \pm 0.0044$. It
also scores 0.11 above \tfunet{}, and it does so while having been trained on
\aoflagger{}'s labels, which themselves score only 0.5698 against the human
expert: the model ends up agreeing with the expert more than its own training
labels do. It is the highest score of all the methods compared in this report:
above Akeret et al.'s \tfunet{} (0.5482), Mesarcik et al.'s U-Net (0.5876) and
RFI-Net (0.5979).

\begin{table}[H]
  \centering
  \caption{Results on the 109 human-labelled LOFAR test images. Values for
    Mesarcik et al.'s U-Net and RFI-Net are as reported in their
    paper.}
  \label{tab:final}
  \begin{tabular}{lcr}
    \toprule
    Method & max \Fone{} & Parameters \\
    \midrule
    $\sigma$-clipping threshold & 0.4103 & --- \\
    \tfunet{} (lr $10^{-4}$, 3 seeds) & $0.5482 \pm 0.0139$ & 465{,}986 \\
    \aoflagger{} & 0.5698 & --- \\
    Mesarcik et al.\ U-Net & $0.5876 \pm 0.0031$ & --- \\
    RFI-Net & 0.5979 & --- \\
    \hybrid{}, class-weighted (3 seeds) & $0.6467 \pm 0.0079$ & 593{,}842 \\
    \textbf{\hybrid{} (3 seeds)} & $\mathbf{0.6603 \pm 0.0040}$ & 593{,}842 \\
    \bottomrule
  \end{tabular}
\end{table}

\begin{figure}[H]
  \centering
  \includegraphics[width=0.8\textwidth]{hybrid_lofar_ladder}
  \caption{Maximum \Fone{} on the 109 LOFAR test images. Error bars show the
    standard deviation over three seeds where available.}
  \label{fig:final}
\end{figure}

A wider version of the model, with 32 filters and 9.3 million parameters,
scored 0.6571 with class weighting (one run), against $0.6467 \pm 0.0079$ for
the narrow class-weighted model. Fifteen times more parameters therefore buys
very little, consistent with the synthetic results in \cref{tab:width}.

\section{Checking the result}
\label{sec:audit}
Because the result is better than the published state of the art, I checked
it for mistakes before relying on it, using a separate script that shares no
evaluation code with the training script.
\begin{itemize}
  \item \textbf{Metrics recomputed independently} from the saved model agree
    with the reported values to within $3\times10^{-5}$; the remaining
    difference comes from GPU non-determinism.
  \item \textbf{No leakage:} none of the training or validation images appears
    in the test set, and the normalisation range is computed from training
    images only.
  \item \textbf{Honest threshold:} choosing the threshold on the validation set
    costs 0.006 in \Fone{} compared with choosing the best threshold on the
    test set itself; the lower, honest value is the one reported.
  \item \textbf{Little overfitting:} maximum \Fone{} is 0.711 on training
    images, 0.682 on validation and 0.659 on test.
  \item \textbf{Not a copy of \aoflagger{}:} the model agrees with \aoflagger{}
    at an \Fone{} of only 0.743. It finds 16{,}485 true RFI pixels that
    \aoflagger{} missed and misses 7{,}345 that \aoflagger{} found.
  \item \textbf{Not a simple frequency rule:} flagging the 16 most
    RFI-affected frequency channels everywhere scores only 0.065.
  \item \textbf{Robust to the choice of test images:} resampling the 109 test
    images 2000 times (bootstrap) gives a 95\% interval for the pooled \Fone{}
    of $[0.615, 0.692]$; RFI-Net's 0.5979 lies outside it, and 99.5\% of the
    resamples beat it.
\end{itemize}
The width of that interval, about $\pm0.04$, is the largest uncertainty in the
result. It comes from having only 109 test images, not from the model.
